In this section, we talk of a new optional feature of Lazart.
The possibility to generate a control flow graph from the normal or mutated version.
%--------------------------------------------------------------CFG--------------------------------------------------------------------------------
\subsection{Control flow graph}
Reminder: A Control Flow Graph is a representation, using graph notation, of all paths that might be traversed through a function during its execution, or control flow.
We generated a control flow graph from the generated LLVM code, which is stored in a ".ll" file.

With Lazart, it is also possible to generate a control flow graph from the mutated code produced from Wolverine or from the program compiled.
We used the control-flow graphs to identify and compare the instructions used for each language (C and Rust) in order to understand their differences.\\

See the example in Appendix A in Figure \ref{fig:annexe1}.

The control flow graph was primarily used to understand the optimization of \texttt{rustc} with their own boolean and give an additional explanation of one of the findings with Table \ref{tbl:results-fissc-vp}: \texttt{verify\_pin} (see below in section \ref{FISSC}) in the section \ref{result section}.